This paper evaluated whether multi-agent LLM architectures can 
improve oracle resolution accuracy over single-model baselines. 
Using 1,189 resolved prediction market questions from KalshiBench, 
we compared two architectures, independent aggregation and 
deliberative consensus, against single-model baselines in a domain 
where correctness has direct financial consequences.

This work makes three contributions to the intersection of multi-agent 
AI systems and decentralized oracle design: (1) empirical evidence 
that independent aggregation offers modest but consistent accuracy 
gains while deliberation degrades performance in evidence-based 
resolution; (2) a characterization of error correlation patterns 
specific to oracle resolution, showing that shared training data and 
architectural lineage produce failures that fundamentally limit 
aggregation benefits; and (3) escalation criteria that use inter-agent 
disagreement and confidence as principled routing signals for hybrid 
AI-human systems. The remainder of this chapter develops each in turn.

The results yield a nuanced answer to the central research question. 
Independent aggregation with confidence-weighted voting achieves the 
highest accuracy at 83.43\%, a statistically significant improvement 
over the best single-model baseline, with the largest gains in 
interpretively ambiguous categories like Politics (+3.2 percentage 
points). Yet these gains fall well short of the theoretical Condorcet 
ceiling of approximately 90\%, a gap explained by moderate-to-high 
error correlations across model pairs ($r = 0.529$--$0.689$): when 
one model fails, the others typically fail on the same question for 
the same underlying reason. The same bound shows up at the category 
level. Well-documented domains like Sports (92.9\%) and Financials 
(93.3\%) approach the ceiling reported by prior single-LLM studies, 
while ambiguous domains like Crypto (66.1\%) and Companies (58.0\%) 
resist improvement from any architecture. Together, these findings 
leave a 14\% hard core of questions that no multi-agent approach 
resolves, motivating principled escalation to human arbitration.

The most striking finding is that deliberative consensus degrades performance rather than improving it. All final-round outputs in Architecture B fall below every single-model baseline, with the best deliberative result trailing the worst individual model by over four percentage points. We attribute this to persuasive error propagation and LLM sycophancy: on genuinely difficult questions where models initially disagree, a confidently wrong model can present compelling but flawed reasoning that causes a correct model to revise its answer. This stands in contrast to prior work on multi-agent debate for arithmetic and logical reasoning tasks, where correct derivations are verifiable step-by-step. Oracle resolution lacks this verifiability, making deliberation vulnerable to confident miscalibration.

These findings carry practical implications for the design of oracle systems as prediction markets continue to grow in volume and complexity. For deployed systems where fast, reliable resolution is the priority, independent aggregation with confidence weighting is both simpler and more accurate than deliberation. The escalation framework provides a structured alternative to ad-hoc human override policies: unanimous, high-confidence resolutions achieve 97.87\% accuracy while covering nearly half of all questions, enabling oracle systems to auto-resolve the easy cases while routing genuinely uncertain questions to human arbitration.

More broadly, these findings suggest that the benefits of multi-agent LLM systems may be fundamentally constrained by correlated reasoning errors across models. In domains where correctness depends on interpreting ambiguous real-world evidence rather than deriving verifiable logical conclusions, debate-based architectures risk amplifying persuasive but incorrect arguments rather than correcting them. As decision-making systems increasingly operate as human-AI hybrids, designing reliable AI oracle systems may require architectures that prioritize independence and calibrated escalation over deliberative consensus.